% !Mode:: "TeX:UTF-8"

%%% 说明: 此部分需要自己填写的内容:  论文创新点.

%%% ``郑重申明'' 和``论文使用授权协议书'' 分别在 statement.tex 和 protocol.tex 中, 无需改动.

%%%%%%%%%%%%%%%%%%%%%%%%%%%%%
%%% -------------  英文封面 (无需改动)-------------   %%%
%%%%%%%%%%%%%%%%%%%%%%%%%%%%%
\thispagestyle{empty}
\renewcommand{\baselinestretch}{1.5}  %下文的行距
\vspace*{0.5cm}

\begin{center}{\zihao{2} \the\Etitle \par}\end{center}

\vfill

\begin{center}
\zihao{4}
\begin{tabular}{ r l }
 Candidate:      &  {\sc \the\Eauthor}      \\
 Supervisor:     &  {\sc \the\Esupervisor}   \\
 Major:          & \the\Emajor  \\
 Speciality:     & \the\Especiality
\end{tabular}

\vspace*{2cm}
\begin{center}
  \iflib % 向图书馆提交电子文档, 使用黑白校徽.
  \includegraphics[height=4cm]{whu.eps}       %%  黑白的. 很小, 只有 10k.
  \else
     \ifprint % 文档打印, 使用黑白校徽.
  \includegraphics[height=4cm]{whu.eps}       %%  黑白的.
  \else
  \includegraphics[height=4cm]{whulogo.eps} %%  彩色的.
  \fi
  \fi
\end{center}


\zihao{-2}
\the\Schoolname\\
{\sc Wuhan University}

\vspace*{1.0cm}

\the\Edate

\end{center}
%%%%%%%--判断是否需要空白页-----------------------------
  \iflib
  \else
  \newpage
  \thispagestyle{empty}
  \cleardoublepage
  \fi
%%%%%%%-------------------------------------------------
{\pagestyle{empty}
%%%--- 加入``郑重声明'' --- %%%%%%%%%%%%%%%%%
% !Mode:: "TeX:UTF-8"
%%-----------------------论文原创性声明-----------------------------------------%
\newpage
\vspace*{20pt}
\begin{center}{\ziju{0.8}\zihao{-2}\heiti 论文原创性声明}\end{center}
\par\vspace*{30pt}
\renewcommand{\baselinestretch}{2}
{\zihao{4} \songti %



本人郑重声明: 所呈交的学位论文, 是本人在导师指导下, 独立进行研究工作所取得的研究成果.
除文中已经标明引用的内容外, 本论文不包含任何其他个人或集体已发表或撰写的研究成果.
对本文的研究做出贡献的个人和集体, 均已在文中以明确方式标明.
本声明的法律结果由本人承担. 

\vskip2cm

\hspace*{4cm}学位论文作者(签名): \hspace{4cm} \hfill \\[1cm]
\hspace*{10cm}年 \hfill  月 \hfill 日\hspace{1cm}\hfill\par}

%%%%%%%--判断是否需要空白页-----------------------------
  \iflib
  \else
  \newpage
  \cleardoublepage
  \fi
%%%%%%%-------------------------------------------------
\renewcommand{\baselinestretch}{1.6}
\small\normalsize
%%%%%%%%%%%%%%%%%
%%% ---加入``武汉大学学位论文使用授权协议书'' ---  %%%%%%
% !Mode:: "TeX:UTF-8"

%%%%%---武汉大学学位论文使用授权协议书---%%%%%%%%%%%%
%%%%%%%%%%%%%%%%%%%%%%%%%%%%%%%%%%%
%%%%%%%%%%%%%%%%%%%%%%%%%%%%%%%%%%%
\newpage\vspace*{20pt}\thispagestyle{empty}
\begin{center}{\zihao{-2}\heiti 武汉大学学位论文使用授权协议书}\end{center}
\par\vspace*{30pt}

本学位论文作者愿意遵守武汉大学关于保存、使用学位论文的管理办法及规定,
即：学校有权保存学位论文的印刷本和电子版, 并提供文献检索与阅览服务；
学校可以采用影印、缩印、数字化或其它复制手段保存论文；
在以教学与科研服务为目的前提下, 学校可以在校园网内公布部分及全部内容.
\begin{enumerate}[1、]
  \item  在本论文提交当年, 同意在校园网内以及中国高等教育文献保障系统(CALIS)、高校学位论文系统提供查询及前十六页浏览服务.
  \item  在本论文提交~$\Box$~当年/~$\Box$~一年/~$\Box$~两年
            /~$\Box$~三年以后，同意在校园网内允许读者在线浏览并下载全文，学校可以为存在馆际合作关系的兄弟高校用户提供文献传递服务和交换服务.(保密论文解密后遵守此规定)
\end{enumerate}

\vskip 15mm

论文作者(签名)：\raisebox{-1ex}{\underline{\makebox[5cm][c]{}}}
\vskip2em
				          				
学\qquad\qquad\quad 号：\raisebox{-1ex}{\underline{\makebox[5cm][c]{}}}
\vskip2em	
					
学\qquad\qquad\quad 院：\raisebox{-1ex}{\underline{\makebox[5cm][c]{}}}					

\vskip  2cm
\begin{flushright}
 日期：\hskip2cm 年\hskip1.2cm 月\hskip1.2cm 日
\end{flushright}

%%%%%%%%%%%%%%%%%%%%%%%%%%%%%%%%%%%%%%%
%%%%%%%--判断是否需要空白页-----------------------------
  \iflib
  \else
  \newpage
  \thispagestyle{empty}
  \cleardoublepage
  \fi
%%%%%%%-------------------------------------------------
%%%%%%%%%%%%%%%%%%
%%%%%%%%%%%%%%%%%%%%%%%%%%%%%%%
}




%%%%%%%%%%%%%%%%%%%%%%%%%%%%%%%
%%% ------------------- 论文创新点----------------------- %%%
%%%%%%%%%%%%%%%%%%%%%%%%%%%%%%%

\newpage\vspace*{20pt}\thispagestyle{empty}
\begin{center}{\zihao{-2}\heiti 论文创新点}\end{center}
\par\vspace*{30pt}
\baselineskip=20pt


  本文的主要创新点在于:

\begin{enumerate}[(1)]
  \item  给出了武汉大学博士毕业论文 \LaTeX{} 模板. 说明了编译方法.

  \item  强调了公式排版的一些细节. 并用例子说明.

  \item  强调了标点符号等的用法. 指出了初学者的一些常见错误.


\end{enumerate}

















%%%%%%%%%%%%%%%%%%%%%%%%%%%%%%%
%%%%%%%--请勿删除以下内容 -------------------------------%
%%%%%%%--判断是否需要空白页-----------------------------%
  \iflib
  \let\cleardoublepage\clearpage
  \else
  \newpage
  \thispagestyle{empty}
  \cleardoublepage
  \fi
%%%%%%%---------------------------------------------------%



